\section*{Chapter \ref{ch:iv2:crypto-for-traders-and-researchers} --- Crypto for Traders and Researchers}
\markboth{\protect\omTocMark{}{Answers: Crypto}}{\protect\omTocMark{}{Answers: Crypto}}

\begin{solution}{iq:iv2:crypto-for-traders-and-researchers:1}
You pay \$100 per interval, \$300 a day and about \$109\,500 a year, which is 10.95\%. When the perpetual is above spot, longs generally pay shorts.

$0.0001 \times \$1$ million $= \$100$, three times a day, 365 days a year.

Funding is what holds a contract with no expiry close to spot. If the perpetual is rich, holding it long costs money, so longs leave and shorts come in.

\iqwrong{Treating 0.01\% as trivial because it looks small. It compounds over 1\,095 intervals a year into a real carry.}
\end{solution}

\begin{solution}{iq:iv2:crypto-for-traders-and-researchers:2}
It is about 10.1\% annualised. Buy spot and short the future, which locks in \$1\,500 per bitcoin, or 2.5\% over the 90 days.

$1\,500/60\,000 = 2.5\%$, and $2.5\% \times 365/90 = 10.14\%$. At expiry, the future converges to spot, so the two legs lock in the difference.

\iqwrong{Comparing the 2.5\% with an annual rate without scaling it. Also, you need capital on both legs: the true return on capital is lower than the basis.}
\end{solution}

\begin{solution}{iq:iv2:crypto-for-traders-and-researchers:3}
You make 3.33 bitcoin at 60\,000 and lose 5 bitcoin at 40\,000. In dollars, that is \$200\,000 made and \$200\,000 lost.

$10^6(1/50\,000 - 1/60\,000) = 3.33$, and $10^6(1/50\,000 - 1/40\,000) = -5$.

The dollar P\&L is the coin P\&L times the current price: $N(P_1/P_0 - 1)$, which is linear. Coin P\&L isn't: you lose more coins in the fall than you gain in the rise.

\iqwrong{Assuming it behaves like a contract that pays dollars. The dollar P\&L does match, but you hold it in a coin whose value has moved, and your margin moves too.}
\end{solution}

\begin{solution}{iq:iv2:crypto-for-traders-and-researchers:4}
At about 45\,250.

The margin is 10\% of notional, or 5\,000 per bitcoin. Liquidation comes when equity falls to the 0.5\% maintenance level, 250 per bitcoin. So the price can fall by 4\,750, to 45\,250.

\iqwrong{Saying 45\,000. That is the bankruptcy price, where the equity is zero. Liquidation has to happen before that.}
\end{solution}

\begin{solution}{iq:iv2:crypto-for-traders-and-researchers:5}
It costs about 22\,222 USDC, an average of 2\,222 per ETH. The marginal price afterwards is about 2\,469.

The invariant is $k = 100 \times 200\,000 = 2 \times 10^7$. After the trade, $x = 90$, so $y = 222\,222$. The marginal price is $y/x = 2\,469$.

\iqwrong{Paying the spot 2\,000 on all 10. The price moves along the curve as you buy, and you pay for that: 11\% on average for 10\% of the pool.}
\end{solution}

\begin{solution}{iq:iv2:crypto-for-traders-and-researchers:6}
Buy at 0.995 and redeem at 0.999, which earns 0.4\% a turn. Delays, redemption limits, or doubts about the reserve can stop the gap from closing. At 1.01, you would mint at 1.00 and sell.

Redeeming gives $1 - 0.001 = 0.999$. Each turn also shrinks the supply, which helps to push the price back.

\iqwrong{Assuming anyone can do this. Usually only approved accounts can mint and redeem directly, so arbitrage depends on them, and a peg can hold or fail on their appetite.}
\end{solution}

\begin{solution}{iq:iv2:crypto-for-traders-and-researchers:7}
No: it is about break-even. The gross gain is \$60, or 10 basis points, and the fees take about \$59.97.

$60\,000 \times 0.0005 + 59\,940 \times 0.0005 = 59.97$. So the net is 3 cents on one bitcoin.

\iqwrong{Comparing prices and stopping. A visible cross has to clear both fees, and it may be gone before both legs fill.}
\end{solution}

\begin{solution}{iq:iv2:crypto-for-traders-and-researchers:8}
Your position is worth about 5.72\% less than holding. It is the same 5.72\% if the price halves.

The pool value is $V = 2\sqrt{kP}$, and holding is worth $(P/P_0 + 1)V_0/2$. So for a price ratio $r$, the ratio is $2\sqrt{r}/(1 + r)$. At $r = 2$ it is $2\sqrt{2}/3 = 0.9428$. At $r = 1/2$ it is the same.

\iqwrong{Thinking the pool loses money in dollars when the price doubles. The pool has gained about 41\%. It is only less than what holding gives, and it comes before fees.}
\end{solution}

\begin{solution}{iq:iv2:crypto-for-traders-and-researchers:9}
You earn about 9.1\% on capital at 20\% margin, and 7.3\% at 50\%. Funding turning negative or a sharp rally that forces you out ends it.

Funding is $0.0001 \times 1\,095 = 10.95\%$ on notional. Capital is the \$1 million of spot and \$200\,000 of margin, so $10.95/1.2 = 9.1\%$. With \$500\,000 of margin, $10.95/1.5 = 7.3\%$.

There's a trade-off: lower margin earns a little more, but a 20\% rally would then wipe out the short's margin, even though the spot leg has gained the same amount on another venue.

\iqwrong{Quoting 10.95\% as the return. That's the return on notional. And the spot gain isn't usable as margin unless the two legs share a venue with cross-margin.}
\end{solution}

\begin{solution}{iq:iv2:crypto-for-traders-and-researchers:10}
Long the future and short the perpetual. Over 90 days you earn 2.7\% in funding and give up 2.5\% of basis, for 0.2\%, or \$2\,000 on \$1 million, if funding holds.

Funding sums to $0.0001 \times 270 = 2.7\%$.

Spot doesn't matter since at expiry the future is worth spot and the perpetual is near spot. The long-future gain is $S - F_0$, and the short-perp gain is $P_0 - S$. Together it is $P_0 - F_0$, which is minus the basis, with no $S$ in it.

The risk is that funding changes. You receive a floating rate and have fixed a basis. So this is a swap: floating funding against a fixed basis. Margin on two legs, and the perpetual trading away from spot at expiry, are the rest.

\iqwrong{Saying this is arbitrage. It would be if funding were fixed. Funding has swung well below this level in calm markets.}
\end{solution}

\begin{solution}{iq:iv2:crypto-for-traders-and-researchers:11}
You are worth \$60\,000 at both. That's a true dollar position. The long's coin P\&L is concave: it gains less in the rise than it loses in the fall.

At 30\,000, the short inverse pays $-60\,000(1/60\,000 - 1/30\,000) = +1$ bitcoin. So you have 2 bitcoin at 30\,000, which is \$60\,000. At 90\,000, it is $-60\,000(1/60\,000 - 1/90\,000) = -1/3$, so $2/3$ of a bitcoin at 90\,000, which is \$60\,000.

For a long with coin margin, a fall both costs you more coins and reduces the value of the coins left. So margin is lost at double speed, and that brings liquidation earlier than a dollar-margined contract.

\iqwrong{Thinking a coin-margined long and a dollar-margined long have the same liquidation risk. For the same dollar notional and margin, the coin version is weaker in a fall.}
\end{solution}

\begin{solution}{iq:iv2:crypto-for-traders-and-researchers:12}
The price ends 12.5\% down. If the liquidations are \$50 million per 1\%, the fall has no limit in this model.

Each 1\% of fall forces \$30 million of selling, and that pushes the price a further $30/50 = 0.6\%$. So the total is $5\% \times (1 + 0.6 + 0.36 + \dots) = 5\%/(1 - 0.6) = 12.5\%$.

At a ratio of 1, the series doesn't converge. So the price falls until the positions or the leverage run out, or the market is halted.

\iqwrong{Treating 5\% as the event. The cascade multiplier, $1/(1 - b)$, is what matters, and $b$ is hidden until leverage is built up.}
\iqfollow{What would you measure in advance?}{Open interest with its estimated liquidation levels, and book depth. A high ratio of the first to the second is a warning.}
\end{solution}

\begin{solution}{iq:iv2:crypto-for-traders-and-researchers:13}
The insurance fund takes the \$10\,000 shortfall. At a fill of 45\,250, the \$5\,000 of remaining margin typically goes to the fund. If the fund is empty, the venue reduces profitable opposite positions, which is called auto-deleveraging.

Below the bankruptcy price of 45\,000, the position has less than zero equity. At 44\,500, that's $500 \times 20 = \$10\,000$. At 45\,250, the equity left is $250 \times 20 = \$5\,000$.

So winning traders can have their hedges cut when they least want it. A market maker who thought it was flat may become outright exposed. That is a real risk to price in.

\iqwrong{Assuming a hedged short can't be hurt by others' liquidations. Auto-deleveraging can close your profitable leg and leave your other leg exposed.}
\end{solution}

\begin{solution}{iq:iv2:crypto-for-traders-and-researchers:14}
It costs about 22\,289 USDC, with a fee of about 67. The band is about 1\,994 to 2\,006.

With the fee on the input, $0.997\,\mathrm{d}y$ goes into the curve. So $\mathrm{d}y = 22\,222/0.997 = 22\,289$, and the fee is 66.87.

A buyer pays $2\,000/0.997 = 2\,006$ at the margin, and a seller receives $2\,000 \times 0.997 = 1\,994$. So inside this band no one gains by trading against the pool.

So the pool is a market maker quoting a fixed half-spread equal to the fee.

\iqwrong{Thinking a fee only matters at the end. It shifts where arbitrage starts, so a pool's price may be off by up to the fee, and arbitrageurs won't correct it.}
\end{solution}

\begin{solution}{iq:iv2:crypto-for-traders-and-researchers:15}
Buy ETH with dollars, sell it for BTC at 0.0340, and sell BTC for dollars. That earns about 2.0\% gross, and about 1.69\% after fees. It would take about 66 basis points per leg to kill it.

The implied rate is $2\,000/60\,000 = 0.03333$, and $0.0340/0.03333 = 1.02$. So \$1 becomes \$1.02 before fees, and $1.02 \times 0.999^3 = 1.0169$ after.

The kill fee solves $(1 - f)^3 = 1/1.02$: $f = 0.66\%$.

\iqwrong{Doing the loop in the wrong direction: $1/1.02$ loses 2\%. Find which leg is rich. Here, ETH is rich against BTC, so you sell ETH for BTC.}
\end{solution}

\begin{solution}{iq:iv2:crypto-for-traders-and-researchers:16}
It's worth about 0.96. A price under 0.92 means the market doubts the other 92\%, or it's a liquidity discount. Size for the case where you can't get out for weeks.

$0.92 + 0.08 \times 0.5 = 0.96$. Even with zero recovery of the stuck 8\%, the coin is worth 0.92. So a price of 0.88 can't come from that 8\% alone.

That price means one of three things: other reserves are also in doubt, redemptions are shut and holders need dollars now, or there's a contagion premium. So the upside is about 9\% if you are right.

\textbf{Sizing.} Assume no redemption for a while, and work out your loss at a much lower price. Also, ask whether you can hold it or need the capital elsewhere.

\iqwrong{Treating the stuck reserve alone as the answer. A run is a problem with timing and access to redemption, as well as with values.}
\end{solution}

\begin{solution}{iq:iv2:crypto-for-traders-and-researchers:17}
About 200 million tokens are minted, the price falls from 1.00 to 0.25, and the average realised is 0.50. So the design turns selling into more selling.

Selling tokens gets dollars out of the pool. To get \$100 million, the dollar side must go from 200 to 100. So the token side goes from 200 to 400, which needs 200 million tokens. The price is $100/400 = 0.25$.

\textbf{The problem.} The backing is the sister token, and its price is what redemption attacks. Each dollar redeemed mints more tokens at a lower price. So once a run begins, more supply is needed for each dollar, in a loop with no floor. Nothing outside the system holds it up.

\iqwrong{Counting 100 million tokens because the price is 1. That ignores your own impact, and the impact doubles the tokens needed.}
\end{solution}

\begin{solution}{iq:iv2:crypto-for-traders-and-researchers:18}
The funding is 0.01\%, 0.01\%, 0.05\%, 0.01\% and $-0.15\%$.

Inside a range for $P$, the clamp cancels the premium, and $F = I$. Beyond it, $F$ follows $P$ with a shift of 0.05\%:
\begin{itemize}
\item $P = 0.02$: $I - P = -0.01$, so $F = 0.01$.
\item $P = 0.06$: $I - P = -0.05$, so $F = 0.01$.
\item $P = 0.10$: $I - P = -0.09$, clamped to $-0.05$, so $F = 0.05$.
\item $P = -0.03$: $I - P = 0.04$, so $F = 0.01$.
\item $P = -0.20$: $I - P = 0.21$, clamped to $+0.05$, so $F = -0.15$.
\end{itemize}

So funding sits at the interest term for $P$ between $-0.04\%$ and $0.06\%$. That's why funding often prints at exactly one value. It acts as a dead zone with a linear slope outside.

\begin{dated}{2026-11}{A published rule}\label{dat:iv2:crypto-for-traders-and-researchers:funding}
One large venue publishes this form, with an interest term of 0.01\% per interval on most contracts, and three intervals a day. Formulas and intervals differ between venues and contracts. Read the current page before you trade.
\end{dated}

\iqwrong{Treating 0.01\% as a market signal of low demand. In the dead zone, it says nothing except that the premium is small.}
\end{solution}

\begin{solution}{iq:iv2:crypto-for-traders-and-researchers:19}
The expected cost is 0.9\% of capital per year. Splitting across three leaves that unchanged and cuts the loss from one failure from 18\% to 6\%, though all three failing would still cost 18\%. A hedged book also loses the hedge.

$0.3 \times 0.05 \times 0.6 = 0.9\%$. With three venues, each with 10\%, the expected cost is $3 \times 0.1 \times 0.05 \times 0.6 = 0.9\%$. A single failure now costs 6\%.

A book long spot on one venue and short perpetual on another is flat only while both exist. When one fails, the other leg is suddenly outright. And the failure tends to come with a violent market, so you may lose much more on that leg than on the frozen funds.

\begin{dated}{2026-11}{One failure's numbers}\label{dat:iv2:crypto-for-traders-and-researchers:fail}
The bankruptcy papers of a large offshore venue that failed in November 2022 show about \$11.2 billion owed to customers on the international platform, against about \$2.5 billion of assets at the petition date.
\end{dated}

\iqwrong{Treating diversification as eliminating venue risk. It reduces the size of each loss, but the expected loss stays the same unless the venues differ in quality.}
\end{solution}

\begin{solution}{iq:iv2:crypto-for-traders-and-researchers:20}
$F = S\kappa/(\kappa - r)$. At $\kappa = 1\,095$ and $r = 5\%$, $F$ is about 0.46 basis points over spot, and the funding is about 5\% a year: it's the interest rate.

\textbf{Derivation.} Try $F = cS$. Hold $c$ of spot, financed at $r$, and short 1 perpetual. Over $\mathrm{d}t$, the spot and the perpetual cancel. Interest costs $rcS\,\mathrm{d}t$, and you receive funding $\kappa(F - S)\,\mathrm{d}t = \kappa(c - 1)S\,\mathrm{d}t$. No risk means no profit: $\kappa(c - 1) = rc$. So $c = \kappa/(\kappa - r)$.

$c - 1 = 0.05/1\,094.95 = 4.57 \times 10^{-5}$. The funding is $\kappa(c - 1) = rc \approx 5\%$.

So funding is the cost of carry, and the premium is what it takes to pay that cost. With a yield $q$ on the coin, it's $\kappa - (r - q)$. As $\kappa \to \infty$, $F \to S$.

\iqwrong{Thinking funding must fall to zero with no premium. A zero premium with a positive rate would give free carry, since a long pays no interest and has full exposure.}
\iqfollow{Why does real funding differ from this?}{Leverage demand and margin costs push it up, and capital between venues doesn't move freely. The formula is the frictionless anchor, and it gives the sign of the excess.}
\end{solution}

\begin{solution}{iq:iv2:crypto-for-traders-and-researchers:21}
About 68\% at 10x, 43\% at 5x and 7.7\% at 2x, using the barrier rule.

With no drift, the chance that a log price touches a barrier $b$ within $T$ is $2\Phi(-b/(\sigma\sqrt{T}))$. Here, $\sigma\sqrt{T} = 0.8\sqrt{30/365} = 0.229$. So $b = \ln 1.1$, $\ln 1.2$ and $\ln 1.5$ give $2\Phi(-0.415) = 0.678$, $2\Phi(-0.795) = 0.427$ and $2\Phi(-1.77) = 0.077$.

A drop of leverage from 5x to 2x reduces the chance about five and a half times. Going from 10x to 5x gives only about 1.6 times. So the risk is highly non-linear in the barrier distance.

\iqwrong{Using $\Phi$ of the terminal distribution, which would halve these numbers. Touching is what matters, as you are liquidated when the barrier is hit, even if the price comes back.}
\iqfollow{What does a simulation give?}{Slightly less at a daily check, since discrete monitoring misses some touches. Real liquidation engines watch continuously, so use the continuous number.}
\end{solution}

\begin{solution}{iq:iv2:crypto-for-traders-and-researchers:22}
You make about 20 basis points of capital a day, or about 73\% a year. Each \$1 million moved has a standard deviation of about \$6\,000. You remove it by hedging the move on a perpetual, which costs fees and a little basis.

Revenue is $\$5 \text{ million} \times 0.001 = \$5\,000$. Transfers cost $\$2 \text{ million} \times 0.0005 = \$1\,000$. So $\$4\,000/\$2 \text{ million} = 0.2\%$ a day.

The risk is $0.8\sqrt{0.5/8\,760} = 0.60\%$ on each \$1 million while the coin is in transit: \$6\,000. That's a day and a half's profit on one \$1 million move if it goes against you at 1 standard deviation. Short a perpetual while the coin is in transit, and unwind it on arrival. That's two more fees, say 5 basis points each: \$1\,000 per \$1 million moved, or \$2\,000 a day. The net falls to \$2\,000, or 10 basis points of capital a day, or about 37\% a year.

The 73\% isn't scalable. An edge like this is competed away. Doubling capital doesn't double \$5 million of opportunity, and as more firms chase it, the 10 basis points shrink.

\iqwrong{Scaling the return to any capital. Arbitrage has limited capacity. And a venue failing would cost more than a year's profit.}
\end{solution}

\begin{solution}{iq:iv2:crypto-for-traders-and-researchers:23}
$V = 2\sqrt{kP}$ has delta $\sqrt{k/P}$, and $\Gamma = -\tfrac12\sqrt{k}\,P^{-3/2}$. The hedged loss is $\tfrac12|\Gamma|\sigma^2P^2\,\mathrm{d}t = V\sigma^2/8\,\mathrm{d}t$. At 80\%, that's 8\% a year, and with value decay, about 7.7\% in expectation over a year.

\textbf{Derivation.} $\tfrac12 \times \tfrac12\sqrt{k}P^{-3/2} \times \sigma^2P^2 = \tfrac14\sqrt{kP}\,\sigma^2 = (V/2)\sigma^2/4 = V\sigma^2/8$.

For a delta-hedged position, $\mathrm{d}V - \Delta\,\mathrm{d}P = \tfrac12\Gamma\sigma^2P^2\,\mathrm{d}t$. So it is pure gamma cost. Since $V$ has drift $-\sigma^2/8$, expected value after a year is $1 - e^{-0.08} = 7.69\%$. A simulation matches this. This loss is called \index{loss versus rebalancing}loss versus rebalancing.

\textbf{Why it doesn't vanish.} The hedge is rebalanced continuously. Each move of the pool sells low and buys high against the arbitrageur, and those losses are locked in. Impermanent loss, measured against holding, can return to zero. This cannot.

\iqwrong{Treating LP fee income as profit. Fees are what you are paid for being short gamma, and you must compare them to this running cost.}
\end{solution}

\begin{solution}{iq:iv2:crypto-for-traders-and-researchers:24}
You lose \$20\,000 at 2\,420. That's exactly the impermanent loss against holding. When the price returns, your P\&L goes back to zero, with no hedge trades, so it was a true path-independent loss.

$k = 2 \times 10^9$. At 2\,420, $V = 2\sqrt{kP} = 2 \times 2.2 \times 10^6 = \$4.4$ million. It has gained \$400\,000. The short loses $1\,000 \times 420 = \$420\,000$. Net: $-\$20\,000$.

Holding would have been $1\,000 \times 2\,420 + 2\,000\,000 = \$4.42$ million. So the impermanent loss is the same \$20\,000, or $2\sqrt{1.21}/2.21 - 1 = -0.45\%$.

So a static hedge at the start gives exactly the impermanent loss: it's path-independent. A continuous hedge gives a realised path-dependent loss of $V\sigma^2/8\,\mathrm{d}t$, which accumulates. For a big one-way move, static loses more. For a volatile round trip, the static loses nothing, and the continuous hedge has paid each wiggle.

That's a gamma-trade choice: static is exposed to the final level, and continuous pays realised variance.

\iqwrong{Thinking impermanent loss and loss versus rebalancing are the same thing. They are equal on average over a large set of paths with the same variance, but they differ on any one path.}
\end{solution}

\begin{solution}{iq:iv2:crypto-for-traders-and-researchers:25}
It multiplies capital by about 20.4. At the top and bottom of the range, it loses about 2.32\% and 2.83\% against holding, and full range loses 0.11\% and 0.14\%. Beyond the range, it holds a single asset and earns no fees.

\textbf{Capital.} For liquidity $L$, full range needs $2L\sqrt{P}$ of value, and the range needs $L(2\sqrt{P} - \sqrt{a} - P/\sqrt{b})$. Their ratio at $a = 1\,800$, $b = 2\,200$, $P = 2\,000$ is 20.4.

\textbf{Loss.} For the range position, compare $V(P_1) = L(2\sqrt{P_1} - \sqrt{a} - P_1/\sqrt{b})$ with what its starting holdings would be worth. At $P_1 = 2\,200$, the ratio is $-2.32\%$. For full range, it's $2\sqrt{1.1}/2.1 - 1 = -0.11\%$. So loss per dollar goes up by the same factor of about 20, as does the rate of loss versus rebalancing and the fee income.

\textbf{Beyond the range.} Above 2\,200, the position is all dollars, and misses any further rise. Below 1\,800, it's all ETH, and takes the full fall. Fees stop in both cases.

So concentration scales everything with leverage. It's like a short strangle at those strikes, and it's not a free increase in yield.

\iqwrong{Saying ``20x capital efficiency means 20x the return.'' It's 20x the fees and 20x the gamma cost, on a position that can leave its range and stop earning.}
\end{solution}

\begin{solution}{iq:iv2:crypto-for-traders-and-researchers:26}
The attacker can buy about 0.47 ETH, and makes about \$216. The victim's extra cost is about \$223. The gap goes to the pool's fees. Batching trades at a single price, or private order flow, would make it harmless.

The victim's quote is 22\,289, so the limit is 22\,512. The attacker buys $a$ first, and the victim's cost rises with $a$. Solve for the cost reaching 22\,512: $a = 0.470$.

The attacker pays 947.60, the victim pays 22\,511.98, and the attacker sells for 1\,163.48. So the profit is 215.88. The extra cost to the victim is $22\,511.98 - 22\,289.09 = 222.89$. And the difference of about 7 goes to fees and LPs.

The profit scales with the tolerance. At a limit of 0.05\%, it is only about \$11, and that may not pay for the transaction. Also, transaction priority costs bite into its margin, and rivals bid it away to validators.

Harmless versions: batch auctions with one clearing price, private transaction routes, or a tighter slippage limit, though that risks failed trades.

\iqwrong{Thinking the slippage limit is a price guarantee. Any tolerance is a gift: it tells the attacker how much can be taken.}
\end{solution}

\begin{solution}{iq:iv2:crypto-for-traders-and-researchers:27}
It costs about 0.0913 bitcoin, which is about \$5\,478. The coin delta is about 0.454, against 0.546 for the dollar option. By parity in coin, $C - P = (S - K)/S = 0$ at the money.

The coin payoff times $S$ is $(S - K)^+$, so it's the same dollar option, and $C_{\$} = S\Phi(d_1) - K\Phi(d_2)$. With $d_1 = -d_2 = 0.115$: $C_{\$} = 5\,478$. So $C_{\text{coin}} = 0.0913$. A rule of thumb: $0.4 \times 0.8 \times \sqrt{30/365} = 0.0917$.

The dollar delta is $\Phi(d_1) = 0.546$. The coin price is $C_{\$}/S$. So the coin-price sensitivity scaled by $S$ is $\Phi(d_1) - C_{\$}/S = 0.454$. The gap is the premium: you've paid in a coin, and that coin itself has delta. To hedge in dollars, you should count the coin premium or margin too.

\iqwrong{Hedging at 0.546 while reading the premium in coins with no adjustment. You would then be over-hedged by the premium's coin exposure, if the premium is held in coin.}
\end{solution}

\begin{solution}{iq:iv2:crypto-for-traders-and-researchers:28}
Net return on capital is about 21\%. With venue risk, it's about 18.5\%, and the Sharpe is about 1.7, not 10. The test is also one year, and it's the wrong year.

\textbf{Costs.} $20 \text{ bp} \times 2 \text{ legs} \times 12 = 4.8\%$ of notional. So $30 - 4.8 = 25.2\%$ on notional. With margin of 20\%, capital is 1.2 times notional, so $25.2/1.2 = 21\%$.

\textbf{Venue.} A 5\% chance of losing half costs 2.5\%, giving a mean of 18.5\%. The standard deviation of that event is $0.5\sqrt{0.05 \times 0.95} = 10.9\%$. With 2\% of other noise, the total is 11.1\%. So $18.5/11.1 = 1.67$.

\textbf{What else.} A one-year test with no losing month contains no event of the kind that kills this trade. A high Sharpe with a hidden tail is the point.

\begin{dated}{2026-11}{How funding moves}\label{dat:iv2:crypto-for-traders-and-researchers:hist}
On one large venue, annualised bitcoin perpetual funding averaged about 17\% in 2020, 31\% in 2021, 4\% in 2022, 8\% in 2023, 12\% in 2024, 5\% in 2025, and 3\% in 2026 through September. In 2022, about a fifth of the prints were negative. So the 30\% test is one year in six, and any one year says little.
\end{dated}

The rest of the list: sample in one hot year; hedging legs assumed perfectly filled; fees at taker level understated; margin calls not modelled; venues that later failed excluded; capacity (funding falls as capital comes in); and a stated Sharpe measured from daily marks in which the tail didn't appear.

\iqwrong{Testing the number and accepting the strategy. Ask what would have to happen for this to lose half, and whether that is in the data. Here it isn't.}
\end{solution}

\begin{solution}{iq:iv2:crypto-for-traders-and-researchers:29}
At 2x, 5x and 10x, gross yield is 6.7\%, 8.3\% and 9.1\%. Monthly liquidation chances are 7.7\%, 42.7\% and 67.8\%. Net yields are about 6.0\%, 4.1\% and 1.7\%. Use low leverage, perhaps about 2x, unless you can move margin fast.

\textbf{Gross.} Capital buys $N$ of spot and $N/L$ of margin, so notional is $N = C/(1 + 1/L)$. At $L = 2$, it's $0.667C$, so $10\% \times 0.667 = 6.67\%$. At 5, $8.33\%$. At 10, $9.09\%$.

\textbf{Liquidation.} The barrier distance is $1/L$ of the short. So $\ln 1.5$, $\ln 1.2$ and $\ln 1.1$, as computed with $2\Phi(-b/(\sigma\sqrt{T}))$ and $T = 30/365$.

\textbf{Net.} Each liquidation costs 1\% of notional. Twelve monthly chances at $p$ gives $12p \times 1\% \times N/C$:
\begin{itemize}
\item 2x: $0.925\% \times 0.667 = 0.62\%$, so 6.05\%.
\item 5x: $5.12\% \times 0.833 = 4.27\%$, so 4.07\%.
\item 10x: $8.13\% \times 0.909 = 7.39\%$, so 1.70\%.
\end{itemize}
So going for leverage to earn 2.4\% more of gross yield costs 6.8\% more in liquidation losses. A cost of 1\% is also on the light side for a violent market, and that pushes further towards low leverage.

\textbf{What else, before running it.}
\begin{itemize}
\item \textbf{Funding regime.} 10\% is a mean. Have a rule to exit when funding stays below your costs. Backtests must include low-funding years.
\item \textbf{Venue limits.} Cap each venue. Know what happens to the other leg if one fails, and pre-plan the exit.
\item \textbf{Margin plumbing.} Cross-margin on one venue is far safer, if you can tolerate the concentration. Otherwise, you need a transfer plan that's faster than the move, and spare capital.
\item \textbf{Auto-deleveraging.} Your profitable short can be cut.
\item \textbf{Stablecoin.} If your dollars are stablecoins, you also hold their peg risk.
\item \textbf{Reporting.} Track carry, fees, hedge slippage and tail losses separately so that you can see when the edge has gone.
\end{itemize}

\iqwrong{Choosing the highest yield. A basis trade looks riskless until the leverage turns a price move into a realised loss, so high leverage on a riskless-looking trade is how such desks die.}
\end{solution}

\begin{solution}{iq:iv2:crypto-for-traders-and-researchers:30}
With no fee, you lose about 7.8\%, or \$780\,000. Arbitrage fee income is about 0.7\%, 2.9\% and 5.2\%. To break even you need yearly uninformed flow of about \$1.4 billion at 5 basis points, \$163 million at 30, and \$26 million at 100. I would take it only if I can show the flow, because high fees drive it away, and arbitrage alone never pays.

\textbf{Simulation.} Hourly prices at 80\% volatility, an arbitrageur moving the pool to the edge of the fee band each hour, and a hedge of the pool's delta. The net is what a delta-hedged LP has without other flow.

\begin{center}
\begin{tabular}{rrrrr}
\toprule
Fee (bp) & Loss & Arb fees & Net & Volume needed (\$m/yr) \\
\midrule
0 & 7.8\% & 0 & $-7.8\%$ & --- \\
5 & 7.8\% & 0.7\% & $-7.1\%$ & 1\,430 \\
30 & 7.8\% & 2.9\% & $-4.9\%$ & 163 \\
100 & 7.8\% & 5.2\% & $-2.6\%$ & 26 \\
\bottomrule
\end{tabular}
\end{center}

The zero-fee loss matches the theory: $1 - e^{-\sigma^2/8} = 7.7\%$.

Arbitrage fees alone are always less than the loss, and they must be: the arbitrageur only trades if it profits, and a hedged LP's result is the negative of that profit. So what's left is what uninformed flow must pay. At 30 basis points, $\$10 \text{ million} \times 4.9\% = \$489\,000$, and dividing by 0.3\% gives \$163 million, or about \$447\,000 a day.

\textbf{Should I take it?} I would look at actual uninformed flow as a share of pool, by fee level. Higher fees need much less flow, but they also lose flow to cheaper pools. The right question is what volume I'd truly get at each fee.

\textbf{Lower vol.} At 40\% volatility, the loss falls to a quarter, about 2\%, since it depends on $\sigma^2$. So it's a volatility-selling position, and I would compare it with selling options directly.

\textbf{Concentrated.} With a 20x range, the fees, the gamma and the loss scale by 20 in the range, and it can go out of range. The break-even flow per unit of liquidity stays about the same, so the flow you need per dollar of capital is about 20 times larger, and flow is shared with other LPs. And the risk of being left outside after a jump is new.

\textbf{Practical.} The hedge costs perpetual funding on a short, it has fees, and it can't be continuous. Each adds to the loss. And smart contract risk, plus a stablecoin leg, isn't in the table.

\iqwrong{Comparing fee income with impermanent loss over a period when the price ended where it began, and concluding that LPs have a free yield. A hedged LP pays variance every day, whatever the final price.}
\end{solution}
